公开卖方观点已经高度一致，预期差不再是“造船会复苏”或“重组会扩大规模”。最新五份合并后原始报告均给出正向评级，2026E归母集中在165.50亿--200.36亿元；13份归档原始报告均为买入或强烈推荐，但实际只有四家券商反复覆盖。下一阶段股价更依赖利润质量与持续期超预期，而不是新的同质化评级。

\section{E13｜当前可比卖方矩阵}

当前横截面仅采用2026Q1后、合并后75.26亿股口径的五份原始PDF。合并前44.7243亿股报告只进入历史修订，不参与当前一致预期。原始PDF均未披露可用目标价，因此本节标题强调“公开研究预期”，而不是声称拥有完整Street目标价。

\begin{exhibitbox}[公开券商预测对照]
\scriptsize
\begin{tabularx}{\exhibitboxwidth}{L{2.0cm}C{1.4cm}L{1.5cm}R{1.8cm}R{1.8cm}R{1.4cm}X}
\toprule
\textbf{券商/来源} & \textbf{日期} & \textbf{评级} & \textbf{2026E收入} & \textbf{2026E归母} & \textbf{EPS} & \textbf{目标/方法} \\
\midrule
华源 BR-01 & 07-14 & 买入 & 1,735.01 & 181.98 & 2.42 & 目标未披露；报告价PE \\
诚通 BR-02 & 05-06 & 强烈推荐 & 1,782.37 & 170.90 & 2.27 & 目标未披露；PE \\
东吴 BR-03 & 05-05 & 买入 & 1,880.07 & 200.36 & 2.66 & 目标未披露；PE \\
华源 BR-04 & 05-01 & 买入 & 1,735.01 & 176.38 & 2.34 & 目标未披露；PE \\
国金 BR-05 & 04-29 & 买入 & 1,924.07 & 165.50 & 2.199 & 目标未披露；PE \\
\rowcolor{accentblue!8} AStock Base & 07-22 & 中性偏多 & 1,775.00 & 172.76 & 2.296 & 2027E周期正常化PE \\
\bottomrule
\end{tabularx}
\sourcenote{BR-01--BR-05均为原始PDF；AStock为独立冻结增长/估值模型。金额单位：亿元。}
\end{exhibitbox}

2026 House归母较Street中位176.38亿元低约2.1\%，分歧不大；2027 House 216.49亿元较中位235.38亿元低约8.0\%；2028 House 246.81亿元较中位299.43亿元低约17.6\%。差异随期限扩大，说明House没有复制卖方对高毛利持续期的假设。

\section{一致看多并不等于十三个独立投票}

13份原始PDF来自四家券商，包含同一团队连续更新。评级一致说明多头叙事已拥挤，却不能按十三个独立观点计数。合并后样本的预测可以作为盈利外部边界，但缺少中性/空头样本使风险分布偏斜。

\begin{exhibitbox}[共识、House与真正分歧]
\small
\begin{tabularx}{\exhibitboxwidth}{L{3.2cm}L{4.0cm}L{4.0cm}X}
\toprule
\textbf{变量} & \textbf{市场/卖方共识} & \textbf{AStock观点} & \textbf{验证方式} \\
\midrule
2026利润方向 & 高价订单集中交付、盈利上修 & 方向可信，不机械H1年化 & 正式半年报与全年交付 \\
毛利持续期 & 高端/绿色与协同继续抬升 & 2026改善，2027--2028不自动外推 & 核心业务毛利与成本 \\
合并协同 & 规模、船位与采购协同自然兑现 & 只对可量化效率与现金给信用 & 费用、周转、少数股东 \\
现金转化 & 高合同负债提供安全性 & 预付款不是自由现金 & OCF、存货、合同资产 \\
资产/军工期权 & 可能构成催化 & 正式公告前全部0 & 交易范围、价格、盈利与股本 \\
目标价 & 市场可能引用弱高目标 & 原始PDF未披露，权重0 & 券商官方原件 \\
\bottomrule
\end{tabularx}
\sourcenote{BR-01--BR-13、CF-03--CF-05；House与variant perception。}
\end{exhibitbox}

这个矩阵的投资含义是：盈利上修本身已经不是稀缺信息。正式H1若只落在预告中点附近、现金和毛利没有额外改善，股价未必因“符合预期”上涨；反之，利润、毛利、现金三项同步超预期才可能形成新的预期差。

\section{弱目标价不能成为Street锚}

华泰58.96元仅来自第三方预览，未取得原始报告；国泰海通47.00元线索对应本地失败页面，没有可用正文。两者均不能核验盈利、倍数、日期、股本与风险披露，因此valuation weight为0。13份原始券商PDF没有可用目标价，也不允许用报告中的PE和EPS反推“目标价”。\srcid{AG-03, AG-04, BR-01--BR-13}

\begin{sourcequalitybox}[目标价证据缺口]
有效Broker/Street目标价为“未披露”，不是“平均约某一价格”。最终目标价由AStock内在价值85\%与可观察市场锚15\%组成，Broker权重严格为0。该缺口降低外部锚的置信度，但不能通过弱来源补齐。
\end{sourcequalitybox}

缺少有效目标价并不阻断自建估值，因为盈利预测、现价与股本仍可复算；它阻断的是“相对Street折价”说法。我们只能比较House与卖方盈利持续期，不能声称目标价高于或低于有效Street平均。

\section{最强反方：等待验证会不会错过上修？}

反方认为，H1预告、在手覆盖和2026交付计划已经足以证明盈利上行，按卖方2027E EPS现价仅约10倍，等待正式现金流会错过估值与盈利双击。这个观点并非稻草人：若高价订单持续至2028、成本受控、现金同步改善，当前Base确实会被快速上修。

但反方成立需要四项同时出现：H1靠近或超过预告上端；核心业务毛利率继续改善；经营现金流与利润同步；2027--2028排程或订单结构获得新的可审计证据。当前只看到利润预告和聚合订单，市场又已高度看多，因此等待验证是对拥挤度和非对称下行的控制，而不是否认周期。

\begin{keyinsight}[预期差的“所以呢”]
真正的variant不是比卖方更看多2026，而是更重视2027--2028利润持续性和现金。动作上，半年报上端利润若没有毛利与现金配合，只能维持；三项同步超预期才升级；任何一项显著低于阈值都应先下修模型。
\end{keyinsight}

本章将公开研究放回正确位置：它是市场已经相信什么的证据，不是报告的声音，更不是目标价替代品。下一章估值将以独立合并盈利为主体，用现价作有限市场锚，并把Broker目标价权重维持为0。
